\section*{Hoofdstuk \ref{ch:g9:climate-past-and-present} --- Het klimaat van de Aarde, vroeger en nu}

\begin{solution}{exo:g9:climate-past-and-present:1}
(a) Ongeveer \qty{-19}{\degreeCelsius}. (b) Ongeveer \qty{15}{\degreeCelsius}.
(c) $15 - (-19) = 34$: ongeveer \qty{34}{\degreeCelsius}. (d) Twee van:
koolstofdioxide, waterdamp, methaan.
\end{solution}

\begin{solution}{exo:g9:climate-past-and-present:2}
(a) \omterm{def:g1:watching-the-weather:weather}{Weer}. (b) \omterm{def:g5:weather-and-climate:climate}{Klimaat}. (c) \omterm{def:g5:weather-and-climate:climate}{Klimaat} (en \omterm{def:g9:climate-past-and-present:climate-change}{klimaatverandering}: een verandering
over decennia). (d) \omterm{def:g1:watching-the-weather:weather}{Weer}.
\end{solution}

\begin{solution}{exo:g9:climate-past-and-present:3}
(a) \qty{19}{cm}, elf jaar. (b) $19 \div 11 \approx \qty{1.7}{cm}$.
(c) De \omterm{def:g4:clouds-rain-snow:precipitation}{sneeuw} van elk jaar valt boven op de \omterm{def:g4:clouds-rain-snow:precipitation}{sneeuw} van de jaren ervoor,
dus hoe dieper een laag, hoe vroeger ze viel.
\end{solution}

\begin{solution}{exo:g9:climate-past-and-present:4}
(a) $22.4 - 21.0 = +\qty{1.4}{\degreeCelsius}$. (b) $20.3 - 21.0 =
\qty{-0.7}{\degreeCelsius}$. (c) \qty{0}{\degreeCelsius}. (d) Die maand
was kouder dan zijn gemiddelde over 1991--2020.
\end{solution}

\begin{solution}{exo:g9:climate-past-and-present:5}
(a) De jaren 1880 tot en met de jaren 1930, de jaren 1950 en de jaren
1960. (b) Ongeveer \qty{0.25}{\degreeCelsius} en
\qty{0.80}{\degreeCelsius}. (c) $0.80 -
0.25 = \qty{0.55}{\degreeCelsius}$. (d) Nee: de decennia van de jaren 1880
tot de jaren 1910 werden iets kouder, en de jaren 1950 waren iets kouder
dan de jaren 1940; maar sinds de jaren 1970 is elk decennium warmer
geweest dan het vorige.
\end{solution}

\begin{solution}{exo:g9:climate-past-and-present:6}
(a) $(324.62 - 315.98) \div 10 \approx \qty{0.86}{ppm}$ per jaar. (b)
$(427.35 - 401.01) \div 10 \approx \qty{2.63}{ppm}$ per jaar. (c) $2.63
\div 0.86 \approx 3.1$: ongeveer drie keer sneller. (d) $(427.35 - 315.98)
\div 315.98 \times 100 \approx 35~\%$.
\end{solution}

\begin{solution}{exo:g9:climate-past-and-present:7}
(a) $8 \times 10^5$ jaar. (b) Ongeveer \num{120000} tot \num{130000} jaar
geleden: de temperatuurkromme komt daar boven 0. (c) Ongeveer
\qty{275}{ppm}. (d) Ja: de pieken en \omterm{def:g1:landscapes-around-us:landform}{dalen} van de twee krommen vallen op
dezelfde momenten.
\end{solution}

\begin{solution}{exo:g9:climate-past-and-present:8}
(a) $3.2 \times 10 = \qty{32}{mm} = \qty{3.2}{cm}$. (b) $2100 - 2025 =
75$ jaar; $75 \times 3.2 = \qty{240}{mm}$, ongeveer \qty{24}{cm} (en meer
als de stijging versnelt). (c) Zeewater zet uit als het opwarmt; ijs op
het land smelt en zijn water stroomt naar de \omterm{def:g2:water-on-earth:ocean}{zee}. (d) Nee: \omterm{def:g8:ice-on-earth:sea-ice}{zee-ijs} drijft
al op de \omterm{def:g2:water-on-earth:ocean}{zee}, dus het smelten ervan doet de \omterm{def:g2:water-on-earth:ocean}{zee} niet stijgen.
\end{solution}

\begin{solution}{exo:g9:climate-past-and-present:9}
(a) 90. (b) 66. (c) 99. (d) \emph{Waarschijnlijk} betekent minstens 66
kansen op 100, geen zekerheid: de krant laat de uitspraak zekerder klinken
dan de wetenschappers zeiden.
\end{solution}

\begin{solution}{exo:g9:climate-past-and-present:10}
Eén koude winter op één plek is \emph{weer}: hij zegt wat er gedurende een
paar maanden op één plek gebeurde. De \omterm{def:g9:climate-past-and-present:global-warming}{mondiale opwarming} is een
verandering van \emph{\omterm{def:g5:weather-and-climate:climate}{klimaat}}, gemeten over decennia, en \emph{mondiaal}:
het is het gemiddelde over de hele Aarde. Een koude winter hier kan
voorkomen terwijl de Aarde als geheel warmer is dan vroeger, en de
gemiddelden per decennium blijven stijgen.
\end{solution}

\begin{solution}{exo:g9:climate-past-and-present:11}
(a) $280 - 190 = \qty{90}{ppm}$, in ongeveer $2 \times 10^4$ jaar. (b) $427
- 280 = \qty{147}{ppm}$, in $2025 - 1750 = 275$ jaar. (c) De tweede
stijging is groter (147 tegen 90~ppm). Tempo: $90 \div (2 \times 10^4) =
4.5 \times 10^{-3}$~ppm per jaar, tegen $147 \div 275 \approx 5.3 \times
10^{-1}$~ppm per jaar: ongeveer 120 keer sneller. (d) Dat het in 2019
hoger was dan op enig moment in minstens twee miljoen jaar (hoge mate van
vertrouwen).
\end{solution}

\begin{solution}{exo:g9:climate-past-and-present:12}
(a) Ze zijn bijna gelijk: 1.07 tegen \qty{1.06}{\degreeCelsius}. De beste
schatting zegt dat menselijke invloed ongeveer de hele waargenomen
opwarming verklaart. (b) Nee: de natuurlijke oorzaken samen veranderden de
temperatuur met hoogstens \qty{0.1}{\degreeCelsius}, tien keer minder dan
de waargenomen opwarming. (c) Minstens 66 kansen op 100 dat de door de mens
veroorzaakte opwarming tussen 0.8 en \qty{1.3}{\degreeCelsius} ligt.
\end{solution}

\begin{solution}{pb:g9:climate-past-and-present:1}
\textbf{1.} Omdat het de referentieperiode is: de anomalieën worden
gemeten vanaf haar eigen gemiddelde.

\textbf{2.} Kouder, met \qty{0.22}{\degreeCelsius}.

\textbf{3.} Warmer, met \qty{0.98}{\degreeCelsius}.

\textbf{4.} $0.98 - (-0.22) = \qty{1.20}{\degreeCelsius}$.

\textbf{5.} $427 - 316 = \qty{111}{ppm}$.

\textbf{6.} $111 \div 316 \times 100 \approx 35~\%$.

\textbf{7.} $427 - 300 = \qty{127}{ppm}$.

\textbf{8.} Koolstofdioxide absorbeert een deel van het infrarode licht dat
het warme oppervlak uitzendt en stuurt een deel ervan terug naar beneden.
Met meer ervan in de lucht komt er meer energie terug naar het oppervlak,
dat daardoor opwarmt.

\textbf{9.} $2025 - 1993 = 32$ jaar; $32 \times 3.2 \approx
\qty{102}{mm}$, ongeveer \qty{10}{cm}.

\textbf{10.} Het opwarmende zeewater zet uit; ijs op het land (\omterm{def:g8:ice-on-earth:glacier}{gletsjers},
\omterm{def:g8:ice-on-earth:ice-sheet}{ijskappen}) smelt en stroomt naar de \omterm{def:g2:water-on-earth:ocean}{zee}.

\textbf{11.} Ja, van dezelfde grootte: ongeveer \qty{1.2}{\degreeCelsius}
tegen \qty{1.09}{\degreeCelsius}. Ze verschillen een beetje omdat de
perioden niet dezelfde zijn (1880--1899 in plaats van 1850--1900;
2015--2024 in plaats van 2011--2020: de recentere jaren waren warmer), en
omdat het IPCC verscheidene archieven combineert, niet alleen dat van de
NASA.

\textbf{12.} ``Sinds het eind van de 19e eeuw is het oppervlak van de
Aarde opgewarmd met ongeveer \textbf{\qty{1.2}{\degreeCelsius}}.''
\end{solution}
